\backchapter{The Notes of the Original}

\noindent The original ends with a section headed ``Notes.'' It is not reproduced here. This account
says what the section contains, which works the book rests on, and what the few notes that go beyond
a reference add to the chapters.

\ssection{What the Notes Are}
The notes are not numbered footnotes. Following the plan set out in the foreword and prologue, which
keeps references out of the reader's way, they are short paragraphs gathered at the back of the book
under chapter headings. Each names the source of a particular description, anecdote, quotation,
table or text box in its chapter, usually with publisher, year and pages. There are about 130 of
them, for thirty-nine of the forty-eight chapters; chapters 22, 26, 27, 29, 30, 31, 34, 35 and 45
have none. Most chapters have between one and five notes. The longest groups belong to chapter 5 and
to the new chapters 47 and 48. A last group, headed ``For Further Study,'' recommends reference
works.

\ssection{The Sources Behind the Chapters}
The notes bear out what the prologue admits, that Shelley borrowed freely to keep the story moving.
He says where a passage is ``reshaped,'' condensed, abridged or adapted from another writer, and
where the quotations of a section all come from one book. A handful of works return again and again.
\emph{Eerdmans' Handbook to the History of Christianity}, edited by Tim Dowley, which he rates among the best first books on the subject, lies behind the portraits of Irenaeus and
Tertullian (chapter 3), the Christological debate (11), the account of icons (15), the summary of
Anabaptist beliefs (25) and the section on Christians under Hitler (42). Paul Hutchinson and Winfred
Garrison's \emph{20 Centuries of Christianity} supplies the contrast between Jewish and Christian
proselytizing (4), the picture of Diocletian and Galerius (9) and the attitude of the East India
Company (38). Roland Bainton's \emph{Christendom} is used for Cyprian on forgiveness (7), for Jerome
(12) and for the baptism of Clovis (16), and his \emph{Here I Stand}, in Shelley's judgment the best
biography of Luther, for the quotations of chapter 24. Williston Walker's \emph{Great Men of the
Christian Church} stands behind Benedict (12), the contrast of Pelagius and Augustine (13) and
Zinzendorf (33), and William Hordern's \emph{A Layman's Guide to Protestant Theology} behind the
description of orthodoxy (5), an anecdote about the Trinity (10) and the spokesmen for liberalism
(40).

For political history Shelley leans on a standard college text, \emph{Civilization Past and Present}
by Wallbank, Taylor and Bailkey. From it come Charlemagne's rise, the explanation of feudalism and
the conflicts of the Germanic kingdoms (18), Urban's call to the crusade (19), the scholar Olbert
(20), the quotation from Holbach (32), the events of the French Revolution (36), and Dickens, Marx
and Engels (41). Other debts are to August Franzen's \emph{A History of the Church} (8, 36),
Friedrich Gontard's \emph{The Chair of Peter} (14, 18), Robert Payne's \emph{The Fathers of the
Western Church} (9, 17), Philip Schaff (10, 17), Henry Bettenson's \emph{Documents of the Christian
Church} for the Old Roman Creed and the Definition of Chalcedon (5, 11), Anne Fremantle's \emph{Age
of Faith} (19, 20), Stephen Neill (15, 38), Herbert Workman (6, 23), Ernest Marshall Howse's
\emph{Saints in Politics}, which Shelley found most helpful for the Clapham chapter (37), and two
articles in \emph{Time} magazine, on Luther (24) and on Loyola and the Jesuits (28). The first two
chapters take their vivid detail from the Reader's Digest volume \emph{Great People of the Bible}, and Bamber Gascoigne's \emph{The Christians} is used in chapters 1 and 9.

Ancient and original writers are cited directly: \emph{The Martyrdom of Polycarp}, Pliny's letter to
Trajan, Eusebius's \emph{Ecclesiastical History} and \emph{Life of Constantine}, Origen's
\emph{Against Celsus} and \emph{On First Principles}, Tertullian's \emph{Prescriptions Against the
Heretics}, Sozomen, Augustine's \emph{Confessions} and \emph{City of God}, Loyola's \emph{Spiritual
Exercises}, Pascal's \emph{Pens\'ees}, the Anabaptist \emph{Martyrs' Mirror} and Rauschenbusch's
\emph{A Theology of the Social Gospel}.

The notes added for this edition name the sources of the new text boxes and chapters. Rodney Stark's
\emph{The Rise of Christianity} lies behind the box in chapter 3. J.~N.~D.~Kelly's \emph{Early
Christian Doctrines}, Craig Evans's \emph{Fabricating Jesus} and a lecture given by Simon Gathercole
in 2012 lie behind the three boxes on Gnosticism in chapter 5. Handbooks of doctrine by Justo
Gonz\'alez, Roger Olson, Allan Coppedge, A.~H.~Strong and Tony Lane are named for the Trinity and
the person of Christ (10, 11), Everett Ferguson's \emph{Church History} for chapter 7, and Alister
McGrath's \emph{Historical Theology} for the theology and the boxes of chapter 24 and for
Renaissance humanism (32). David Bentley Hart's \emph{The Story of Christianity} (18, 21), Robert
Wilken (19, 32), Glenn Sunshine on the Inquisition (23), Thomas Kidd's \emph{The Great Awakening}
(37), James Leo Garrett's \emph{Baptist Theology} on Andrew Fuller (38) and Garry Wills's \emph{Under
God} on the Scopes trial (43) are cited in the same way, and in chapter 38 the reviser points to an
essay of his own, ``The Hermeneutics of Conversion,'' on the immediacy with which Scripture is read. For chapter 46 the notes name
George Weigel's \emph{The Cube and the Cathedral}, Hart's \emph{Atheist Delusions} and McGrath's
\emph{The Twilight of Atheism} on the ``smarter'' atheism, and Tom Doyle's \emph{Dreams and Visions}
on stirrings among Muslims. For chapters 47 and 48 they name Philip Jenkins, Mark Noll's \emph{The
New Shape of World Christianity}, Allan Anderson on Pentecostalism, the Pew Forum's report
\emph{Global Christianity} of 2011, David Aikman's \emph{Jesus in Beijing}, and Noll and Carolyn
Nystrom's \emph{Clouds of Witnesses}.

\ssection{Remarks That Add to the Text}
A few notes say more than where a passage comes from. In chapter 18 Shelley remarks that the politics
of the Holy Roman Empire can confuse the layman, which is why he followed a textbook there.

One note to chapter 47 enlarges on the Pentecostals, who stood ``in the vanguard of the Southern
Counter-Reformation.'' The movement itself dates from the first years of the 1900s and arose mainly in North America. A century later it counted at least 400 million people, heavily concentrated in the
Global South, and by about 2040 there could be as many as a billion, far more than the world's
Buddhists and roughly as many as its Hindus. Other notes explain the chapter's figures. The estimate
that charismatics are about one in four of all Christians is the Pew Forum's, and the table of
Pentecostal numbers is adapted from Mark Noll, who relies on David Barrett's \emph{World Christian
Encyclopedia}.

Two notes printed under chapter 47 belong with the account of the emerging churches in chapter 46.
D.~A.~Carson describes the heart of the movement, which some of its leaders prefer to call a
``conversation,'' as the belief that a changed culture calls for a new kind of church. Leaders are expected to adapt; those who refuse, the argument runs, do not see how far inherited forms have hidden the gospel from a younger generation.
Many observers, the note adds, have associated the emerging church with postmodernity. David Wells
is then quoted on the postmodern outlook. It takes so many shapes that there are only postmodern
perspectives, but all of them hold that meaning has died, above all the rational meaning the
Enlightenment provided, and they find no grounds for any other kind.

A note to chapter 48 adds that South Korea's churches were still gaining almost 6 percent a year after 2000, thanks to energetic evangelism and to a government that valued Christianity as a counterweight to communism. Another says that the chapter's ``conservative numbers'' are those of the Pew
Forum, and a last one recommends the online \emph{Dictionary of African Christian Biography}.

\ssection{For Further Study}
The section closes with two paragraphs of advice for readers who want more than the reading lists at
the end of the chapters and the notes supply. For any topic in the book they recommend two reference
works with articles and further reading: \emph{The Oxford Dictionary of the Christian Church}, third
edition, of which F.~L.~Cross and E.~A.~Livingstone are the editors, and J.~D.~Douglas's \emph{The New International Dictionary of the Christian Church}. For the chapters on American
Christianity they add the \emph{Dictionary of Christianity in America}, edited by Daniel G.~Reid, and
for maps Franklin H.~Littell's \emph{The Macmillan Atlas History of Christianity}.

\ssection{The Reading Lists}
Every chapter of the original also ends with a short list of suggested books, more than three hundred
entries in all. About sixty of them carry the asterisk that marks the titles newly recommended by the
reviser. Those lists, like the notes, are left to the original.
